\subsubsection*{Orientation}

\begin{remark*}[Working vocabulary]
This orientation uses truth assignments, logical consequence, tautology,
logical equivalence, formula contexts, and counterexample assignments to justify
common argument forms.
\end{remark*}

\begin{remark*}[Exposition]
Argument forms are reusable patterns of inference. The purpose of this section
is to justify the basic inference tools that will be used throughout the rest of
the project. The justification is semantic: an argument form is valid when no
truth assignment makes every premise true while making the conclusion false.
\end{remark*}

\begin{remark*}[Section map]
This section introduces argument forms, substitution instances, validity,
invalidity, and counterexample assignments. It then proves the validity of the
standard inference patterns, isolates common invalid forms, explains rules of
replacement, and connects informal proof patterns with semantic consequence.
Formal proof systems are previewed but not developed here.
\end{remark*}

\begin{remark*}[Boundary]
The goal is rigorous senior-undergraduate propositional logic. The section
proves that the proof moves used later are truth-preserving, but it does not
build Hilbert systems, sequent calculi, cut elimination, or metatheoretic
completeness machinery.
\end{remark*}
